\documentclass[12pt]{article}
% Préambule commun à tous les documents extraits de « La Revue CAPES »
\usepackage[T1]{fontenc}
\usepackage[utf8]{inputenc}
\usepackage{lmodern}
\usepackage{amsmath,amssymb,amsthm,amsfonts,mathrsfs}
\usepackage{setspace}
\usepackage{listings}
\usepackage{pgfplots}
\pgfplotsset{compat=1.18}
\usepackage{longtable}
\usepackage{adjustbox}
\usepackage{lettrine}
\usepackage{tkz-tab}
\usepackage{changepage}
\usepackage{pdflscape}
\usepackage{graphicx}
\usepackage{tikz}
\usepackage{xcolor}
\usepackage[a4paper,top=2cm,bottom=2.5cm,left=2.5cm,right=2cm]{geometry}
\usepackage{fancyhdr}
\usepackage{draftwatermark}
\SetWatermarkText{La RevueCapes}
\SetWatermarkLightness{0.9}
\SetWatermarkScale{2}
\usepackage{hyperref}
\hypersetup{colorlinks=true,linkcolor=blue!50!black,urlcolor=blue!50!black}

\definecolor{revue}{RGB}{20,60,120}

\pagestyle{fancy}
\fancyhf{}
\renewcommand{\headrulewidth}{0.4pt}
\setlength{\headheight}{15pt}

% \entete{Matière}{Titre}{Type de document}
\newcommand{\entete}[3]{%
  \fancyhead[L]{\small\textsc{La Revue CAPES} -- #1}%
  \fancyhead[R]{\small #2}%
  \fancyfoot[C]{\thepage}%
  \thispagestyle{plain}%
  \begin{center}
    {\color{revue}\rule{\textwidth}{1.2pt}}\\[4pt]
    {\small\textsc{Concours directs CAP-CEG / CAPES -- Burkina Faso}}\\[6pt]
    {\LARGE\bfseries\color{revue} #2}\\[6pt]
    {\large\itshape #1 -- #3}\\[2pt]
    {\color{revue}\rule{\textwidth}{1.2pt}}
  \end{center}
  \vspace{0.5cm}%
}

% Encadré pour les remarques ajoutées dans les corrigés
\newcommand{\remarque}[1]{\par\medskip\noindent\fcolorbox{revue}{revue!6}{\parbox{\dimexpr\linewidth-2\fboxsep-2\fboxrule}{\textbf{\color{revue}Remarque.} #1}}\par\medskip}


\begin{document}
\entete{Mathématiques}{Corrigé détaillé du sujet type n°12}{Correction}
\begin{spacing}{1.5}

\textbf{Exercice 1}


1. Calcul du volume du domaine $D$

Le domaine $D$ est défini par :
\[
D = \{ (x, y, z) \in \mathbb{R}^3 : 0 \leq x \leq 1, 0 \leq y \leq (1 - \sqrt{x})^2, z^2 \leq 2xy \}.
\]

Étape 1 : Description des bornes
\begin{itemize}
    \item $x$ varie de $0$ à $1$,
    \item pour un $x$ fixé, $y$ varie entre $0$ et $(1 - \sqrt{x})^2$,
    \item pour un $(x, y)$ fixé, $z$ varie entre $-\sqrt{2xy}$ et $\sqrt{2xy}$ (d'après $z^2 \leq 2xy$).
\end{itemize}

Étape 2 : Calcul du volume

Le volume est donné par :
\[
V = \int_{0}^{1} \int_{0}^{(1-\sqrt{x})^2} \int_{-\sqrt{2xy}}^{\sqrt{2xy}} 1 \, dz \, dy \, dx.
\]

Intégrale sur $z$ :
\[
\int_{-\sqrt{2xy}}^{\sqrt{2xy}} 1 \, dz = \left[ z \right]_{-\sqrt{2xy}}^{\sqrt{2xy}} = 2\sqrt{2xy}.
\]

Intégrale sur $y$ :
\[
\int_{0}^{(1-\sqrt{x})^2} 2\sqrt{2xy} \, dy = 2\sqrt{2x} \int_{0}^{(1-\sqrt{x})^2} \sqrt{y} \, dy.
\]
Avec $\int \sqrt{y} \, dy = \frac{2}{3} y^{3/2}$, on obtient :
\[
\int_{0}^{(1-\sqrt{x})^2} \sqrt{y} \, dy = \frac{2}{3} \left[ y^{3/2} \right]_{0}^{(1-\sqrt{x})^2} = \frac{2}{3} \left((1-\sqrt{x})^3 \right).
\]
Ainsi :
\[
\int_{0}^{(1-\sqrt{x})^2} 2\sqrt{2xy} \, dy = \frac{4\sqrt{2x}}{3} (1-\sqrt{x})^3.
\]

Intégrale sur $x$ :
\[
V = \int_{0}^{1} \frac{4\sqrt{2x}}{3} (1-\sqrt{x})^3 \, dx.
\]
En posant $u = \sqrt{x}$, $x = u^2$, et $dx = 2u \, du$, les bornes deviennent $u \in [0, 1]$. On obtient :
\[
V = \frac{4}{3} \int_{0}^{1} 2\sqrt{2} u^2 (1-u)^3 \, (2u \, du).
\]
Cela donne :
\[
V = \frac{8\sqrt{2}}{3} \int_{0}^{1} u^3 (1-u)^3 \, du.
\]

Développez $(1-u)^3$, puis intégrez terme par terme pour obtenir le résultat final.



2. Calcul de $T$ avec changement de variables

Le domaine $D$ est défini par :
\[
D = \{ (x, y) \in \mathbb{R}^2 : x + y \leq 1, x \geq 0, y \geq 0 \}.
\]

L'intégrale à calculer est :
\[
T = \int \int_D (x + y)^2 e^{x^2 - y^2} \, dx \, dy.
\]

Étape 1 : Changement de variables
Posons :
\[
x = \frac{u+v}{2}, \quad y = \frac{u-v}{2}.
\]
Le jacobien est donné par :
\[
\frac{\partial (x, y)}{\partial (u, v)} = 
\begin{vmatrix}
\frac{\partial x}{\partial u} & \frac{\partial x}{\partial v} \\
\frac{\partial y}{\partial u} & \frac{\partial y}{\partial v}
\end{vmatrix} =
\begin{vmatrix}
\frac{1}{2} & \frac{1}{2} \\
\frac{1}{2} & -\frac{1}{2}
\end{vmatrix} = -\frac{1}{2}.
\]
Donc, $|J| = \frac{1}{2}$.

Les nouvelles bornes pour $u$ et $v$ sont :
\[
u = x + y \quad \text{et} \quad v = x - y, \quad \text{avec } u \in [0, 1] \text{ et } v \in [-u, u].
\]

Étape 2 : Réécriture de l'intégrale
Avec ce changement, on a :
\[
x + y = u, \quad x^2 - y^2 = uv.
\]
Ainsi :
\[
T = \int_{0}^{1} \int_{-u}^{u} u^2 e^{uv} \cdot \frac{1}{2} \, dv \, du.
\]

Intégrons d'abord sur $v$ :
\[
\int_{-u}^{u} e^{uv} \, dv = \frac{1}{u} \left[ e^{uv} \right]_{-u}^{u} = \frac{1}{u} \left(e^{u^2} - e^{-u^2}\right).
\]

Ainsi :
\[
T = \int_{0}^{1} u^2 \cdot \frac{1}{2u} \left(e^{u^2} - e^{-u^2}\right) \, du.
\]
\[
T = \frac{1}{2} \int_{0}^{1} u \left(e^{u^2} - e^{-u^2}\right) \, du.
\]

Posons $w = u^2$, $dw = 2u \, du$, et $u \in [0, 1]$ implique $w \in [0, 1]$. On obtient :
\[
T = \frac{1}{4} \int_{0}^{1} (e^w - e^{-w}) \, dw.=\frac{1}{4}\left[e^w + e^{-w}\right]_0^1=\frac{1}{4}(e+e^{-1}-2)\]

3. Calcul de $T$ sur le domaine $D$

Le domaine $D$ est donné par :
\[
D = \{ (x, y, z) \in \mathbb{R}^3 : 0 \leq x \leq 1, 0 \leq y \leq 1, 0 \leq z \leq xy \}.
\]

L'intégrale à calculer est :
\[
T = \int_{0}^{1} \int_{0}^{1} \int_{0}^{xy} x^\alpha y^\beta z \, dz \, dy \, dx.
\]

Étape 1 : Intégrale sur $z$
\[
\int_{0}^{xy} z \, dz = \left[ \frac{z^2}{2} \right]_{0}^{xy} = \frac{(xy)^2}{2}.
\]

Étape 2 : Intégrale sur $y$
\[
T = \int_{0}^{1} \int_{0}^{1} x^\alpha y^\beta \cdot \frac{(xy)^2}{2} \, dy \, dx = \frac{1}{2} \int_{0}^{1} \int_{0}^{1} x^\alpha y^{\beta+2} x^2 \, dy \, dx.
\]
\[
T = \frac{1}{2} \int_{0}^{1} x^{\alpha+2} \int_{0}^{1} y^{\beta+2} \, dy \, dx.
\]

Intégrale sur $y$ :
\[
\int_{0}^{1} y^{\beta+2} \, dy = \frac{1}{\beta+3}.
\]

Ainsi :
\[
T = \frac{1}{2(\beta+3)} \int_{0}^{1} x^{\alpha+2} \, dx.
\]

Intégrale sur $x$ :
\[
\int_{0}^{1} x^{\alpha+2} \, dx = \frac{1}{\alpha+3}.
\]

Ainsi, on obtient finalement :
\[
T = \frac{1}{2(\beta+3)(\alpha+3)}.
\]



\medskip\textbf{Exercice 2}

Soit $f$ un endomorphisme nilpotent de $E$ d'indice de nilpotence $p\in  \mathbb{N}$. Soit $x$ un vecteur de $E$ vérifiant $f^{p-1}(x)= 0$.
 
 Montrons tout d'abord la propriété suivante : pour tout $i\in \mathbb{N}$, $f^i(x)$ est différent de 0 si et seulement si $i\leqslant p-1$. En effet, $p$ étant l'indice de nilpotence de $f$, on a $f^p = 0$ et aussi $f^{p+k} =f^k\circ f^p = f^k\circ 0_{L(E)} = 0_{L(E)}$ pour tout $k\in \mathbb{N}^\ast$ (ici $L(E)$ désigne l'ensemble des applications linéaires de $E$ dans $E$, dans la suite on notera tout simplement 0 à la
 place de $0_L(E,E)$). D'autre part, si, par l'absurde, on avait $f^i(x) = 0$ avec $i < p-1$ (le cas
 $i = p-1$ étant impossible par hypothèse sur $x$) il existerait $k\in \mathbb{N}^\ast$ tel que $i + k = p-1$ et $f^{p-1}(x) = (f^k\circ f^i)(x) = f^k(f^i(x)) = f^k(0) = 0$ ce qui est une contradiction avec
 $f^{p-1}(x)\neq 0$.
 
 Montrons maintenant que la famille $(x,f(x),...,f^{p-1}(x))$ est libre. Soit $(\lambda_0,...,\lambda_{p-1})\in \mathbb{R}^p$ tel que
 \[\sum_{i=0}^{p-1}\lambda_if^i(x)=0 \text{  } (\ast)\]
 
 Avec la convention $f^0=Id_{L(E)}$. On compose alors la somme précédente avec $f^{p-1}$
 
 \[f^{p-1}(\sum_{i=0}^{p-1}\lambda_if^i(x))=\sum_{i=0}^{p-1}\lambda_if^{p-1+i}(x)=\lambda_0f^{p-1}(x)\].
 
  Ainsi $\lambda_0f^{p-1}(x) = 0$ et puisque $f^{p-1}(x)= 0$ on en déduit que $\lambda_0 = 0$. On continue alors en procédant par récurrence. Nous venons de faire l'initialisation. Soit $k < p-1$ un
 indice quelconque mais fixé. Supposons que $\lambda_i = 0$ pour tout $i\in \{0,...,k\}$. Montrons que $\lambda_{k+1} = 0$. On compose $(\ast)$ avec $f^{p-2-k}$ :
 
 \[f^{p-2-k}(\sum_{i=k+1}^{p-1}\lambda_if^i(x))=\lambda_{k+1}f^{p-1}(x)=0\]
 
  Ainsi $\lambda_{k+1} = 0$ ce qui conclut l'hérédité et finalement on obtient que $(\lambda_0,...,\lambda_{p-1}) = 0_{\mathbb{R}^p}$
 et la famille $(x,f(x),...,f^{p-1}(x))$ est libre.
 
 b) $E$ étant un espace vectoriel de dimension $n$ on ne peut pas avoir de famille libre possédant $p$ éléments avec $p > n$. Ainsi $p\leqslant n$. Il existe donc $k\in \mathbb{N}$ tel que $p+k = n$ et,
 comme montré précédemment, $f^n = f^{k+p} = f^k\circ f^p = f^k\circ 0 = 0$.
 
 c) Si l'indice de nilpotence de $f$ est $n$ alors, il existe $x\in E$ tel que $f^{n-1}(x)= 0$.
 
 D'après le point a) la famille $(x,f(x),...,f^{n-1}(x))$ est libre. Puisque cette famille possède $n$ éléments il s'agit d'une famille libre maximale dans $E$ (qui est de dimension $n$) et c'est donc une base.
 
 2) Soit $(e_1,...,e_n)$ une base de $E$. Soit $(p_1,...,p_n)\in \mathbb{N}^n$ tel que $f^{p_i}(e_i) = 0$ pour tout $i\in \{1,...,n\}$ (les $p_i$ existent par hypothèse). On note $p_m = max_{1\leqslant i\leqslant n}p_i$ (et on remarque que $f^{p_m}(e_i)=0$ pour tout $i\in \{1,...,n\})$. Soit $x=\sum_{i=1}^nx_ie_i$ un vecteur de $E$,on a:
 
 \[
 f^{p_m}(x)=\sum_{i=1}^nx_if^{p_m}(e_i)=0
 \]
 
 Ainsi $f^{p_m}=0$ et $f$ est un endomorphisme nilpotent. En utilisant 1)b)on en déduit que
 $f^n=0$.
 
 \medskip\textbf{Exercice 3}
 
 
 (1) Explicitons les évènements $(X = 2), (X = 3), (X = 4)$, et déterminons la valeur de $p_2,
 p_3, p_4$.
 
$X=2$ :

Pour que $X=2$, il faut obtenir deux piles consécutifs dès les deux premiers lancers (pile au premier lancer et pile au deuxième lancer).
 
 $P_2=P(X=2)=\frac{2}{3}.\frac{2}{3}=\frac{4}{9}$
 
 \medskip
 
 $X=3$
 
 Pour que $X=3$,  il faut obtenir face au premier lancer, puis pile au deuxième et au troisième lancers:
 
 $P_3=P(X=3)=\frac{1}{3}.\frac{2}{3}.\frac{2}{3}=\frac{4}{27}$.
 
 \medskip
 $X=4$
 
 Pour que $X=4$, il faut obtenir soit (pile, face, pile, pile), soit (face,face, pile, pile) au cours des quatre premiers lancers:
 
 $P_4=P(X=4)=\frac{2}{3}.\frac{1}{3}.\frac{2}{3}.\frac{2}{3}+\frac{1}{3}.\frac{1}{3}.\frac{2}{3}.\frac{2}{3}=\frac{8}{81}+\frac{4}{81}=\frac{12}{81}=\frac{4}{27}$
 


(2) Montrons que l'on a $p_n = \frac{2}{9}p_{n-2}+\frac{1}{3}p_{n-1}, n\geq 4.$

Raisonnons par récurrence :

-Pour $n=4$, on a $P_4=\frac{12}{81}$

$\frac{2}{9}p_{2}+\frac{1}{3}p_{3}=\frac{2}{9}.\frac{4}{9}+\frac{1}{3}.\frac{4}{27}=\frac{8}{81}+\frac{4}{81}=\frac{12}{81}=p_4$ alors l'expression est vrai au rang 4. 

Vérifions le cas $n$

 Si le premier lancer est \( F \), alors \( X - 1 \) suit la même loi que \( X \), donc :
   \[
   P(\text{commence par } F) = \frac{1}{3} p_{n-1}.
   \]
 Si le premier lancer est \( P \), suivi de \( F \), alors \( X - 2 \) suit la même loi que \( X \), donc :
   \[
   P(\text{commence par } PF) = \frac{2}{3} \cdot \frac{1}{3} \cdot p_{n-2} = \frac{2}{9} p_{n-2}.
   \]

En sommant ces deux cas, on obtient la récurrence :
\[
p_n = \frac{2}{9} p_{n-2} + \frac{1}{3} p_{n-1}, \quad \text{pour } n \geq 4.
\]

 
 (3) Déduisons l'expression de $p_n$ pour tout $n$.
 
 La suite $p_n$ est une suite récurrente linéaire du second ordre. 
 
 \[p_n = \frac{2}{9}p_{n-2}+\frac{1}{3}p_{n-1}\Leftrightarrow  p_n -\frac{1}{3}p_{n-1}-\frac{2}{9}p_{n-2}=0\]
 
 L'équation caractéristiques associée est :
 
 $r^2-\frac{1}{3}r-\frac{2}{9}=0$ 
 
 \[\Delta=1\Rightarrow r_1=-\frac{1}{3} \text{ et } r_2=\frac{2}{3}\]

Ainsi la solution générale est :

\[p_n=A\left(\frac{2}{3}\right)^n+B\left(\frac{-1}{3}\right)^n\]

Utilisons les conditions initiales pour déterminer $A$ et $B$.

Pour $n=2$, $p_2=\frac{4}{9}=A\left(\frac{2}{3}\right)^2+B\left(\frac{-1}{3}\right)^2=\frac{4A}{9}+\frac{B}{9}=\frac{4}{9}\Rightarrow 4A+B=4$ $(*)$

Pour $n=3$, $p_3=\frac{4}{27}=A\left(\frac{2}{3}\right)^3+B\left(\frac{-1}{3}\right)^3=\frac{8A}{27}-\frac{B}{27}=\frac{4}{27}\Rightarrow 8A-B=4$ $(**)$
 
 $(*)+(**)\Longrightarrow 12A=8\Longrightarrow A=\frac{2}{3}$
 
$(*)\Longrightarrow B=4-4A=4-4.\frac{2}{3}=\frac{4}{3}$

Ainsi $p_n=\frac{2}{3}\left(\frac{2}{3}\right)^n+\frac{4}{3}\left(\frac{-1}{3}\right)^n$


 
 (4) Rappelons, pour $q \in ]-1,1[$, l'expression de       
$\sum_{n=0}^{+\infty}nq^n$

\[\sum_{n=0}^{+\infty}nq^n=\frac{q}{(1-q)^2} \text{ pour } \vert q\vert<1\]


\quad Calculons alors $E(X)$.
 
 \begin{align*}
 E(X)&=\sum_{n=1}^{+\infty}np_n\\
 &=\sum_{n=1}^{+\infty}n\left(\frac{2}{3}\left(\frac{2}{3}\right)^n+\frac{4}{3}\left(\frac{-1}{3}\right)^n\right)\\
 &=\frac{2}{3}\sum_{n=1}^{+\infty}n\left(\frac{2}{3}\right)^n + \frac{4}{3}\sum_{n=1}^{+\infty}n\left(\frac{-1}{3}\right)^n\\
 &=\frac{2}{3}\times \frac{\frac{2}{3}}{(1-\frac{2}{3})^2}+\frac{4}{3} \frac{\frac{-1}{3}}{(1-\frac{-1}{3})^2}\\
 &=4-\frac{3}{4}\\
 E(X)=\frac{13}{4}
 \end{align*}
 
 
 
 
 
 \medskip\textbf{Exercice 4}
 
 Calculer l'intégrale $I = \int_{0}^{a}\frac{1}{1+f(x)}dx$
 
 \textbf{ Transformation de l'intégrale avec le changement de variable $t = a-x$}

En effectuant le changement de variable $t = a-x$, nous obtenons :

\[
I = \int_0^a \frac{1}{1+f(a-x)}dx.
\]

D'après la relation donnée dans l'énoncé :

\[
f(x) f(a-x) = 1,
\]

nous avons :

\[
f(a-x) = \frac{1}{f(x)}.
\]

Ainsi, en substituant dans l'intégrale :

\[
I = \int_0^a \frac{1}{1+\frac{1}{f(x)}}dx.
\]

\textbf{ Simplification de l'intégrande}

Nous simplifions l'expression dans le dénominateur :

\[
1 + \frac{1}{f(x)} = \frac{f(x)+1}{f(x)}.
\]

Donc :

\[
\frac{1}{1+\frac{1}{f(x)}} = \frac{f(x)}{f(x)+1}.
\]

Ainsi, l'intégrale devient :

\[
I = \int_0^a \frac{f(x)}{f(x)+1}dx.
\]

En remarquant que :

\[
\frac{f(x)}{f(x)+1} = 1 - \frac{1}{f(x)+1},
\]

nous pouvons décomposer l'intégrale :

\[
I = \int_0^a \left( 1 - \frac{1}{f(x)+1} \right) dx.
\]

Ce qui donne :

\[
I = \int_0^a dx - \int_0^a \frac{1}{f(x)+1}dx.
\]

\textbf{ Utilisation de la symétrie}

Par définition de $I$, on reconnaît dans le deuxième terme à droite la même intégrale que $I$ :

\[
I = a - I.
\]

Ce qui donne l'équation :

\[
2I = a.
\]

Ainsi :

\[
I = \frac{a}{2}.
\]


 

\end{spacing}
\end{document}
